% Generated draft: requires mathematical and source review.
\documentclass[11pt]{article}
\usepackage[margin=25mm]{geometry}
\usepackage{fontspec}
\setmainfont{Latin Modern Roman}
\usepackage{amsmath,amssymb,mathtools,unicode-math}
\setmathfont{Latin Modern Math}
\usepackage{enumitem,needspace,xurl,hyperref}
\setlength{\parindent}{0pt}
\setlength{\parskip}{5pt}
\setlength{\emergencystretch}{3em}
\newcommand{\N}{\mathbb{N}}
\newcommand{\Z}{\mathbb{Z}}
\newcommand{\Q}{\mathbb{Q}}
\newcommand{\R}{\mathbb{R}}
\newcommand{\C}{\mathbb{C}}
\newcommand{\F}{\mathbb{F}}
\newcommand{\A}{\mathbb{A}}
\newcommand{\PP}{\mathbb{P}}
\newcommand{\E}{\mathbb{E}}
\newcommand{\eps}{\varepsilon}
\newcommand{\dd}{\,\mathrm{d}}
\begin{document}
\section*{UnsolvedMath ID: 3800001}
\textbf{Source problem code:} catalogue entry
\subsection{Definitions and mathematical statement}
Let $P\subset\R^3$ be a three-dimensional convex polytope. An edge unfolding cuts its boundary along selected polytope edges, keeps the remaining face adjacencies, and develops the resulting connected surface into $\R^2$, isometrically on every face. The usual non-overlap requirement is that the relative interiors of distinct developed faces be disjoint. The source says ``non-self-intersecting'' and does not specify whether additional boundary contacts are forbidden.

The supplied item asks three questions.
\begin{enumerate}
\item Does every such $P$ admit a non-self-intersecting edge unfolding? This is identified in the research summary as Dürer's conjecture, also called Shephard's problem.
\item Form the dual edge graph $G^*(P)$, whose vertices are faces of $P$ and whose edges join faces sharing a polytope edge. Given a spanning tree $T$ of $G^*(P)$, keep the face adjacencies in $T$ and cut the other adjacencies to obtain the corresponding tree development. Does a minimum spanning tree for the source's ``natural dihedral-angle weighting'' always give a non-self-intersecting unfolding? The source does not define that weight function, including interior versus exterior angle or whether the weights increase or decrease with the angle. It also does not specify whether every minimum tree or at least one minimum tree must work when there are ties. Thus the weighted-tree assertion remains a schema until those data are fixed.
\item Does every convex polytope have a creased unfolding that cannot be refolded into a different convex polytope? Here ``different'' would ordinarily mean noncongruent, but the permitted refolding operations are not supplied. In particular the role of prescribed creases, boundary identifications, preservation of the intrinsic surface metric, and allowable new folds must be specified. If $\mathcal R(U)$ denotes the resulting, presently undefined class of convex refoldings of a creased unfolding $U$ of $P$, the intended uniqueness target has the form
\[
 \forall P\ \exists U\quad
 \forall Q\in\mathcal R(U),\qquad Q\text{ is congruent to }P.
\]
This schema does not select a missing refolding convention. The source uses ``every convex polytope'' in this subquestion without explicitly repeating the three-dimensional restriction.
\end{enumerate}
The main edge-unfolding question concerns the given geometry of $P$. Allowing an affine change of its geometry is a separate, weaker version.
\subsection{Short English statement}
The item asks for non-overlapping edge nets of convex polyhedra, whether a dihedral-angle minimum spanning tree always supplies one, and whether a creased net can force a unique convex refolding.

The supplied research summary calls the general edge-unfolding problem open, while reporting Ghomi's affine version solved: every convex polyhedron has a simple edge unfolding after an affine transformation. It interprets this as the absence of a combinatorial obstruction, not a solution for the original fixed geometry. It also reports O'Rourke's result that every polar zonohedron edge-unfolds to a net, and mentions positive classes such as prisms and certain zonohedra. It calls the weighted-dual-tree and refolding questions apparently open. Its final account of those subquestions is truncated, so no missing evidence is reconstructed.

These are source-reported assessments. The source index explicitly warns that its pages have not been systematically updated since 2001; access on 2026-07-29 does not itself establish current status.
\subsection{Sources}
\begin{itemize}
\item Local UnsolvedMath registry, permanent ID 3800001, source code catalogue entry, ``Unfolding convex polytopes,'' attributed to Jeff Erickson and cited proposers. Dataset curation and metadata: Ulam AI and UnsolvedMath Contributors, CC BY 4.0.
\item Erickson, \emph{Open Problems in Computational Geometry}, item ``Unfolding convex polytopes,'' \url{https://jeffe.cs.illinois.edu/open/}; HTML extraction, accessed 2026-07-29. The source warns of no systematic updating since 2001.
\item M. Ghomi, ``Affine unfoldings of convex polyhedra,'' \emph{Geometry \& Topology} 18 (2014), 3055--3090, DOI 10.2140/gt.2014.18.3055, as cited in the supplied summary.
\item J. O'Rourke, ``Polar zonohedra edge-unfold to nets'' (2023), arXiv:2302.07747, as cited in the supplied summary.
\item Source provenance marks status evidence NEEDS\_REVIEW and rights NEEDS\_REVIEW: public access was verified, but redistribution permission was not established. This source-specific qualification is separate from the dataset licence.
\end{itemize}
\end{document}
